\documentclass[a4paper,11pt]{ctexart}

\usepackage[top=1.8cm,bottom=1.8cm,left=2.1cm,right=2.1cm]{geometry}
\usepackage{amsmath,amssymb}
\usepackage{enumitem}
\usepackage{fancyhdr}
\usepackage{lastpage}

\pagestyle{fancy}
\fancyhf{}
\renewcommand{\headrulewidth}{0pt}
\fancyfoot[C]{\small 第~\thepage~页\quad 共~\pageref{LastPage}~页}

\linespread{1.28}
\setlength{\parindent}{0pt}

\newcommand{\tr}{\operatorname{tr}}
\newcommand{\rk}{\mathrm{r}}
\newcommand{\T}{\mathrm{T}}
\newcommand{\blank}[1][2.8cm]{\underline{\hspace{#1}}}

\begin{document}

\begin{center}
  {\Large\bfseries 华中科技大学 2014～2015 学年第二学期}\\[0.4em]
  {\Large\bfseries 《线性代数》课程考试试卷（A 卷）}
\end{center}

\vspace{0.8em}

\textbf{一、判断题（$2\text{ 分}\times 8 = 16\text{ 分}$）（判断题为计算机阅卷，请使用 2B 铅笔规范填涂，填涂区域要丰满、不要使用划线、打勾、打叉等错误填涂方式）}

\begin{enumerate}[label=\arabic*., leftmargin=2em, itemsep=0.65em, topsep=0.5em]
  \item 设 $A$ 为 $n$ 阶方阵（$n \geqslant 3$），则对于任意数 $k$，有 $(kA)^* = kA^*$，$A^*$ 为 $A$ 的伴随矩阵.

  \item 设 $A$ 为 $m \times n$ 阶实矩阵，若 $A^{\T}A = 0$，则 $A = 0$.

  \item 设 $n$ 阶方阵 $A$ 与 $B$ 相似，则对于任意的多项式 $f(\lambda)$，有 $f(A)$ 相似于 $f(B)$.

  \item 若 $n$ 阶行列式 $D_n$ 的数值相同的元素大于 $n^2 - n$ 个，则行列式的值等于零.

  \item 若齐次线性方程组 $AX = 0$ 只有零解，则对应的非齐次线性方程组 $AX = b$ 有唯一解.

  \item 若 $n$ 阶方阵 $A$ 的特征值为 $1$ 或 $0$，则 $A^2 = A$.

  \item 设 $A$ 为 $n$ 阶可逆矩阵，则 $A$ 的任何一行都不是其余行的线性组合.

  \item 设 $A = (a_{ij})_{n \times n}$ 为 $n$ 阶正定矩阵，则 $A$ 的主对角线元素 $a_{ii} > 0$，$i = 1, 2, \cdots, n$.
\end{enumerate}

\vspace{0.9em}

\textbf{二、填空题（$4\text{ 分}\times 5 = 20\text{ 分}$）（请将答案填写在答题卡指定位置上）}

\begin{enumerate}[label=\arabic*., leftmargin=2em, itemsep=1.0em, topsep=0.5em]
  \item $D_4 = \begin{vmatrix}
      1 & x & y & z \\
      x & 1 & 0 & 0 \\
      y & 0 & 1 & 0 \\
      z & 0 & 0 & 1
    \end{vmatrix} = \blank$.

  \item 设矩阵 $A = (a_1, a_2, a_3, a_4)$，其中 $a_1, a_2, a_3, a_4$ 都是四维列向量，$P$ 为四阶矩阵，若 $AP = (a_1, a_2, a_3, a_1 + a_4)$，则 $P = \blank$.

  \item 若 $4$ 阶矩阵 $A$ 的四个列向量 $a_1, a_2, a_3, a_4$ 满足 $a_1 - a_2 + a_3 - a_4 = 0$，且 $\rk(A) = 3$，则线性方程组 $AX = 0$ 的基础解系为 $\blank$.

  \item 已知三阶方阵 $A$ 的特征值为 $1, -1, 2$，设 $B = A^3 - 5A^2$，则 $|B| = \blank$.

  \item 若二次型 $f(x_1, x_2, x_3) = x_1^2 + 4x_2^2 + 4x_3^2 + 2t x_1 x_2 + 2 x_1 x_3 + 4 x_2 x_3$ 是正定的，则 $t$ 的取值范围是 $\blank$.
\end{enumerate}

\vspace{0.9em}

\textbf{三、（8 分）}计算 $n$ 阶行列式
\[
  D_n = \begin{vmatrix}
    a + b & ab & & & \\
    1 & a + b & ab & & \\
    & 1 & a + b & \cdots & \\
    & & \cdots & \cdots & ab \\
    & & & 1 & a + b
  \end{vmatrix}.
\]

\vspace{0.6em}

\textbf{四、（10 分）}已知 $3$ 阶矩阵 $A$ 的逆矩阵为
\[
  A^{-1} = \begin{bmatrix}
    1 & 1 & 1 \\
    1 & 2 & 1 \\
    1 & 1 & 3
  \end{bmatrix},
\]
求 $A$ 的伴随矩阵 $A^*$ 的逆矩阵 $(A^*)^{-1}$.

\vspace{0.6em}

\textbf{五、（12 分）}已知线性方程组
\[
  \begin{cases}
    x_1 + x_2 + x_3 + x_4 = 0, \\
    x_2 + 2x_3 + 2x_4 = 1, \\
    -x_2 + (a - 3)x_3 - 2x_4 = b, \\
    3x_1 + 2x_2 + x_3 + a x_4 = -1,
  \end{cases}
\]
求参数 $a, b$ 的取值，使得方程组
\begin{enumerate}[label=\arabic*., leftmargin=2.4em, itemsep=0.25em, topsep=0.3em]
  \item 有唯一解；
  \item 无解；
  \item 有无穷多解，并求通解.
\end{enumerate}

\vspace{0.6em}

\textbf{六、（12 分）}设矩阵 $A = \begin{bmatrix} 1 & -1 & 1 \\ a & 4 & b \\ -3 & -3 & 5 \end{bmatrix}$ 有三个线性无关的特征向量，$\lambda = 2$ 是 $A$ 的二重特征值，试确定参数 $a, b$ 及可逆矩阵 $P$，使得 $P^{-1}AP$ 为对角矩阵.

\vspace{0.6em}

\textbf{七、（12 分）}用正交变换 $X = CY$ 化二次型 $f(x_1, x_2, x_3) = x_1^2 + x_2^2 + x_3^2 + 4x_1 x_2 + 4x_1 x_3 + 4x_2 x_3$ 为标准型.

\vspace{0.6em}

\textbf{八、（10 分）}设 $A, B$ 为 $n$ 阶方阵，$I$ 为单位矩阵，且 $A^2 = B^2 = I$，$|A| = -|B|$，证明矩阵 $(A + B)$ 不可逆.

\end{document}
